\documentclass[10pt,a4paper,twoside]{article}\usepackage[]{graphicx}\usepackage[]{xcolor}
% maxwidth is the original width if it is less than linewidth
% otherwise use linewidth (to make sure the graphics do not exceed the margin)
\makeatletter
\def\maxwidth{ %
  \ifdim\Gin@nat@width>\linewidth
    \linewidth
  \else
    \Gin@nat@width
  \fi
}
\makeatother

\definecolor{fgcolor}{rgb}{0.345, 0.345, 0.345}
\newcommand{\hlnum}[1]{\textcolor[rgb]{0.686,0.059,0.569}{#1}}%
\newcommand{\hlsng}[1]{\textcolor[rgb]{0.192,0.494,0.8}{#1}}%
\newcommand{\hlcom}[1]{\textcolor[rgb]{0.678,0.584,0.686}{\textit{#1}}}%
\newcommand{\hlopt}[1]{\textcolor[rgb]{0,0,0}{#1}}%
\newcommand{\hldef}[1]{\textcolor[rgb]{0.345,0.345,0.345}{#1}}%
\newcommand{\hlkwa}[1]{\textcolor[rgb]{0.161,0.373,0.58}{\textbf{#1}}}%
\newcommand{\hlkwb}[1]{\textcolor[rgb]{0.69,0.353,0.396}{#1}}%
\newcommand{\hlkwc}[1]{\textcolor[rgb]{0.333,0.667,0.333}{#1}}%
\newcommand{\hlkwd}[1]{\textcolor[rgb]{0.737,0.353,0.396}{\textbf{#1}}}%
\let\hlipl\hlkwb

\usepackage{framed}
\makeatletter
\newenvironment{kframe}{%
 \def\at@end@of@kframe{}%
 \ifinner\ifhmode%
  \def\at@end@of@kframe{\end{minipage}}%
  \begin{minipage}{\columnwidth}%
 \fi\fi%
 \def\FrameCommand##1{\hskip\@totalleftmargin \hskip-\fboxsep
 \colorbox{shadecolor}{##1}\hskip-\fboxsep
     % There is no \\@totalrightmargin, so:
     \hskip-\linewidth \hskip-\@totalleftmargin \hskip\columnwidth}%
 \MakeFramed {\advance\hsize-\width
   \@totalleftmargin\z@ \linewidth\hsize
   \@setminipage}}%
 {\par\unskip\endMakeFramed%
 \at@end@of@kframe}
\makeatother

\definecolor{shadecolor}{rgb}{.97, .97, .97}
\definecolor{messagecolor}{rgb}{0, 0, 0}
\definecolor{warningcolor}{rgb}{1, 0, 1}
\definecolor{errorcolor}{rgb}{1, 0, 0}
\newenvironment{knitrout}{}{} % an empty environment to be redefined in TeX

\usepackage{alltt}
\linespread{1.1}
%paquetes
\usepackage[utf8]{inputenc}
\usepackage{amsfonts}
\usepackage[nottoc,numbib]{tocbibind}
\usepackage{hyperref}
\usepackage[hang,footnotesize,bf]{caption}
\usepackage{amsmath}
\usepackage[left=2.5cm,top=3cm,bottom=3cm,right=2.5cm]{geometry}
\usepackage{booktabs}
\usepackage{authblk}
\usepackage[margin=10pt,font=small,labelfont=bf]{caption}
\usepackage[normalem]{ulem}
\usepackage{bbold}
\usepackage{hyperref}
\usepackage[super]{natbib}
\usepackage{enumitem}
\usepackage{float}


%new_commands
\newcommand{\R}{\textsf{R}}
\newcommand{\Rpackage}[1]{\textsf{#1}}
\newcommand{\Robject}[1]{\texttt{#1}}

% Las librerías que hacen falta por ahora





\IfFileExists{upquote.sty}{\usepackage{upquote}}{}
\begin{document}

\begin{titlepage}
  \begin{center}
    \vspace*{4cm}
    \textbf{\Huge Statistics in Health Sciences}
    \vspace{4em}\\
    \textbf{\huge \underline{Assignment 1}} \\[1em]
    \large{Cluster analysis with missing data using the \texttt{R} packages \texttt{mice} and \texttt{miclust}}\\
    [1em]
    Jordi Ye, NIU: 1751973\\
    Xavier Gaseni Piñol, NIU: 1753108\\
    Alejandro Cabezas Bamba, NIU: 1749626\\
    Xinyu Lou, NIU: 1635130
    \vfill
  \end{center}
\end{titlepage}

\tableofcontents

\newpage

\section{Introduction}

The aim of this assignment is to apply the framework proposed by Basagaña et al. to the data set \texttt{copd} in the file \texttt{copd\_redux.RData}. Data \texttt{copd} is a subset of 45 variables of the original data analyzed in the cited paper.

\section{Analysis}
\subsection{Description of missingness}
Provide the same missingness description than in the original analysis. Specifically, reproduce the following text that is shown in the original paper:\\[10pt]
The required packages are:
\begin{itemize}
  \item Number of participants;
  \item Number of variables;
  \item Missing data percentage of each variable;
  \item Missing data percentage of less than $5\%$ of missing;
  \item Percentage of variables with all variables without any missing;
  \item Percentage among all cell with missing data.
\end{itemize}
\begin{knitrout}\footnotesize
\definecolor{shadecolor}{rgb}{0.969, 0.969, 0.969}\color{fgcolor}\begin{kframe}
\begin{alltt}
\hldef{> }\hlkwd{load}\hldef{(}\hlsng{"copd_redux.RData"}\hldef{)}
\hldef{> }\hldef{n} \hlkwb{<-} \hlkwd{length}\hldef{(copd}\hlopt{$}\hldef{id)} \hlcom{#participant number}
\hldef{> }\hldef{v} \hlkwb{<-} \hlkwd{ncol}\hldef{(copd)} \hlcom{#variables number}
\hldef{> }\hldef{M} \hlkwb{<-} \hldef{n}\hlopt{*}\hldef{v} \hlcom{#cell number}
\hldef{> }
\hldef{> }\hldef{missing} \hlkwb{<-} \hlkwd{colSums}\hldef{(}\hlkwd{is.na}\hldef{(copd))}
\hldef{> }\hldef{miss_percentatge} \hlkwb{<-} \hldef{(missing}\hlopt{/}\hldef{n)}\hlopt{*}\hlnum{100} \hlcom{#missing }
\hldef{> }
\hldef{> }\hldef{minimum_missing_pp} \hlkwb{<-} \hlkwd{min}\hldef{(miss_percentatge)} \hlcom{#minimum missing percentage}
\hldef{> }\hldef{maximum_missing_pp} \hlkwb{<-} \hlkwd{max}\hldef{(miss_percentatge)} \hlcom{#maximum missing percentage}
\hldef{> }
\hldef{> }\hldef{miss_under_5} \hlkwb{<-} \hlkwd{sum}\hldef{(miss_percentatge} \hlopt{<} \hlnum{5}\hldef{)}
\hldef{> }\hldef{pctg_under_5} \hlkwb{<-} \hldef{(miss_under_5}\hlopt{/}\hldef{v)} \hlopt{*} \hlnum{100}
\hldef{> }\hlcom{#percentage of variables with less than 5% missing}
\hldef{> }
\hldef{> }\hldef{complete} \hlkwb{<-} \hlkwd{sum}\hldef{(miss_percentatge} \hlopt{==} \hlnum{0}\hldef{)}
\hldef{> }\hldef{complete_pctg} \hlkwb{<-} \hldef{(complete}\hlopt{/}\hldef{v)}\hlopt{*}\hlnum{100}
\hldef{> }\hlcom{#percentage of patients with complete data on all variables}
\hldef{> }
\hldef{> }\hldef{overall} \hlkwb{<-} \hlkwd{sum}\hldef{(}\hlkwd{is.na}\hldef{(copd))}
\hldef{> }\hldef{overall_pctg} \hlkwb{<-} \hldef{(overall}\hlopt{/}\hldef{M)}\hlopt{*}\hlnum{100}
\hldef{> }\hlcom{#Overall percentage of the information is missing}
\end{alltt}
\end{kframe}
\end{knitrout}
Now we complete the missing values of the text:\\

“The subset used here included \textcolor{red}{342} participants and \textcolor{red}{45} variables. The range of missing data in the different variables went from \textcolor{red}{45}\%, with \textcolor{red}{0}\% to \textcolor{red}{31.6}\%, with \textcolor{red}{68.9}\% of the variables having less than 5\% of the data missing. Only \textcolor{red}{11.1} of the participants had complete data on all  \textcolor{red}{45} variables. Overall, \textcolor{red}{11.1} of the information was missing; that is, \textcolor{red}{5.4}\% of  \textcolor{red}{15,390} cells ($\textcolor{red}{45} \times \textcolor{red}{342} = \textcolor{red}{15,390}$) had missing data.”

\subsection{Multiple imputation}
Apply multiple imputation to \texttt{copd} data using 100 imputations. Explain in a reasoned manner how you decided to set the arguments \texttt{method} and \texttt{predictorMatrix} when using the \texttt{mice()} function.\\[5pt]

Before staring to decide the arguments for \texttt{method} and \texttt{predictorMatrix} we should do a little summary to know the classes of each variable of this \texttt{copd} data-set. 
\begin{knitrout}\footnotesize
\definecolor{shadecolor}{rgb}{0.969, 0.969, 0.969}\color{fgcolor}\begin{kframe}
\begin{alltt}
\hldef{> }\hlkwd{str}\hldef{(copd)}
\end{alltt}
\begin{verbatim}
## 'data.frame':	342 obs. of  45 variables:
##  $ id               : num  3045 5155 5062 5146 1055 ...
##  $ index_barthel2v1 : num  100 100 100 100 100 90 100 100 100 100 ...
##  $ quo_calcpostv1   : num  40.5 55.2 28.6 67 40.5 ...
##  $ rv_prepv1        : num  180 134 202 125 166 ...
##  $ incp_fevv1       : num  10.619 -0.694 -4.255 -1.408 20.976 ...
##  $ fevp_postv1      : num  40 76.9 20.1 70.4 40.7 ...
##  $ rv_tlcv1         : num  54.8 46.5 71.8 46.2 51.8 ...
##  $ ctnf_s_mv1       : num  2.957 0.045 1.023 0.64 0.557 ...
##  $ vtdvdbid_ecov1   : num  30 32 NA 29 NA NA NA NA NA NA ...
##  $ pp_ecov1         : num  11 8 10 10 9.5 9 10 11 9 9 ...
##  $ sin11_v1         : Factor w/ 2 levels "0","1": 2 1 1 2 1 2 1 2 1 1 ...
##  $ vas2_v1          : num  8 4 0 4 5 ...
##  $ ic_prepv1        : num  67.9 96.7 38.7 80.2 85.2 ...
##  $ iam_chv1         : Factor w/ 2 levels "0","1": 1 1 1 1 1 1 1 2 1 1 ...
##  $ sin15_v1         : Factor w/ 2 levels "0","1": 2 1 1 1 2 1 1 1 2 1 ...
##  $ cil8_s_mv1       : num  3.59 2.75 3.72 2.8 6.04 ...
##  $ sin7_v1          : Factor w/ 2 levels "0","1": 1 1 1 2 1 1 1 1 1 1 ...
##  $ sin8_v1          : Factor w/ 2 levels "0","1": 2 2 2 2 2 2 1 1 1 2 ...
##  $ ch_v1            : num  0.4 1 2.6 1.8 0.8 ...
##  $ vc_prepv1        : num  73.4 92.4 49 88.7 73.4 ...
##  $ dtdve_ecov1      : num  63 45 55 60 53 48 49 51 57 49 ...
##  $ cccs1_v0         : Factor w/ 2 levels "0","1": 2 1 1 1 2 1 1 1 1 2 ...
##  $ kcals_set_v0     : num  1635 15055 2265 6054 10965 ...
##  $ impactsv1        : num  29.46 3.55 35.51 24.72 23.22 ...
##  $ monoA_v1         : num  0.41 0.495 0.827 0.543 0.289 ...
##  $ lvmass_ase_ecov1 : num  302914 142108 219991 287863 199797 ...
##  $ vtit_catv1       : Factor w/ 2 levels "0","1": 1 NA 1 2 2 NA NA 1 1 NA ...
##  $ urg_ambv0        : Factor w/ 2 levels "0","1": 1 1 1 1 2 1 1 1 2 1 ...
##  $ po2_v1           : num  63.1 80.5 63.5 92 80 ...
##  $ ahpf_v1          : num  2.5 3.2 3 2.9 4 ...
##  $ pro_v1           : num  6.73 6.93 7.7 6.8 6.9 ...
##  $ basA_v1          : num  0.1104 0.0228 0.0165 0.032 0.0346 ...
##  $ fc_pmv1          : num  0 2 0 1 0 1 5 0 0 0 ...
##  $ tad_v1           : num  80 93 80 90 75 70 94 90 80 86 ...
##  $ col_v1           : num  174 282 NA 174 145 135 168 213 211 212 ...
##  $ sin5_v1          : Factor w/ 2 levels "0","1": 2 2 1 1 2 1 2 1 1 1 ...
##  $ bc_v1            : num  0 0 0 0 0 0 0 0 0 0 ...
##  $ fvcp_postv1      : num  66.1 102.9 43.2 83.5 67.2 ...
##  $ vas1_v1          : num  1 2 1.7 3.2 0 ...
##  $ sin6_v1          : Factor w/ 2 levels "0","1": 2 1 2 2 1 1 1 1 2 2 ...
##  $ graudisn_v1      : Factor w/ 6 levels "0","1","2","3",..: 4 3 3 4 3 6 6 3 3 1 ...
##  $ cpcr_mg_l_mediav1: num  0.08 3.11 1.07 0.08 0.08 ...
##  $ neutA_v1         : num  4.78 6.05 5.22 3.51 2.74 ...
##  $ ahpl_v1          : num  1 2.2 2.3 5.2 1 ...
##  $ activityv1       : num  61.3 35.5 41.8 35.6 36.3 ...
\end{verbatim}
\end{kframe}
\end{knitrout}
As we see, most of the variables are numerical, however there are some variables which are factors. Their names and corresponding numbers of levels are listed below:
\begin{knitrout}\footnotesize
\definecolor{shadecolor}{rgb}{0.969, 0.969, 0.969}\color{fgcolor}\begin{kframe}
\begin{alltt}
\hldef{> }\hldef{categoricas} \hlkwb{<-} \hlkwd{names}\hldef{(copd)[}\hlkwd{sapply}\hldef{(copd, is.factor)]}
\hldef{> }\hldef{categoricas_latex} \hlkwb{<-} \hlkwd{gsub}\hldef{(}\hlsng{"_"}\hldef{,} \hlsng{"\textbackslash{}\textbackslash{}\textbackslash{}\textbackslash{}_"}\hldef{, categoricas)}
\hldef{> }\hldef{niveles} \hlkwb{<-} \hlkwd{sapply}\hldef{(copd[categoricas], nlevels)}
\hldef{> }\hldef{binarias} \hlkwb{<-} \hlkwd{names}\hldef{(niveles[niveles} \hlopt{==} \hlnum{2}\hldef{])}
\hldef{> }\hldef{binarias_latex} \hlkwb{<-} \hlkwd{gsub}\hldef{(}\hlsng{"_"}\hldef{,} \hlsng{"\textbackslash{}\textbackslash{}\textbackslash{}\textbackslash{}_"}\hldef{, binarias)}
\hldef{> }\hldef{no_binarias} \hlkwb{<-} \hlkwd{names}\hldef{(niveles[niveles} \hlopt{>} \hlnum{2}\hldef{])}
\end{alltt}
\end{kframe}
\end{knitrout}
As we can see, this data set has 11 variables that are
factors, and their names are \texttt{sin11\_v1, iam\_chv1, sin15\_v1, sin7\_v1, sin8\_v1, cccs1\_v0, vtit\_catv1, urg\_ambv0, sin5\_v1, sin6\_v1, graudisn\_v1}. The following factors are binary distributed: \texttt{sin11\_v1, iam\_chv1, sin15\_v1, sin7\_v1, sin8\_v1, cccs1\_v0, vtit\_catv1, urg\_ambv0, sin5\_v1, sin6\_v1}



\end{document}
